%titolo{Le foglie}
%autore{}
%album{}
%tonalita{DO}
%famiglia{Scout}
%gruppo{}
%momenti{}
%identificatore{le_foglie}
%data_revisione{2026_07_1}
%trascrittore{Lunardelli Luca}

\beginsong{Le foglie}

\beginverse
\[C]Chi la vede \[G]mai \[F] una foglia \[C]sola se il cielo è \[F]grande 
\[G]e gli alberi son \[C]alti se il baco \[G]già pia\[F]no piano la in\[C]goia 
chi si ac\[G]corgerà, a cosa ser\[F]virà una fo\[C]glia \[G7]sola.
\endverse

\beginchorus
\[C]Il respiro del \[F]bosco e \[C]tante foglie vi\[G]cine 
\[C]e se canta ogni fo\[F]glia \[C]una canzone si sen\[G]te. 
\[F]È così che il vento noi possiamo can\[C]tare 
se noi tutte in\[F]sieme \[C]unite siamo \[F]veramente, \[F]veramen\[C]te. 
\endchorus

\beginverse
\chordsoff
Chi la vede mai una foglia sola 
scende la pioggia 
e piano già se ne va via 
col vento sai se se ne vola 
chi si accorgerà, a cosa servirà 
una foglia sola. Il respiro del bosco e tante foglie\dots{} 
Chi la vede mai una foglia sola se quando è verde 
su quel ramo si riposa se gialla ormai e il sole la indora chi si accorgerà, 
a cosa servirà una foglia sola. Il respiro del bosco e tante foglie\dots{} 
\endverse
\endsong
